\documentclass[11pt]{article}
\usepackage{fontspec}
\usepackage{xeCJK}
\usepackage{xurl}
\usepackage[a4paper,margin=25mm]{geometry}
\usepackage{amsmath,amssymb,amsthm,mathtools}
\usepackage{xcolor}
\usepackage{booktabs,array,enumitem,fancyhdr,hyperref}
\hypersetup{colorlinks=true,linkcolor=blue!50!black,urlcolor=blue!50!black}
\setlength{\emergencystretch}{2em}
\setlength{\parskip}{0.3em}
\linespread{1.16}
\allowdisplaybreaks
\newtheorem{theorem}{定理}[section]
\newtheorem{proposition}[theorem]{命题}
\newtheorem{lemma}[theorem]{引理}
\theoremstyle{definition}
\newtheorem{definition}[theorem]{定义}
\theoremstyle{remark}
\newtheorem{remark}[theorem]{说明}
\newcommand{\dd}{\,\mathrm d}
\newcommand{\ii}{\mathrm i}
\newcommand{\Rea}{\operatorname{Re}}
\newcommand{\Ima}{\operatorname{Im}}
\newcommand{\Li}{\operatorname{Li}}
\newcommand{\supp}{\operatorname{supp}}
\pagestyle{fancy}
\fancyhf{}
\fancyhead[L]{\small 学习笔记}
\fancyhead[R]{\small 初稿}
\fancyfoot[C]{\thepage}
\setlength{\headheight}{15pt}
\author{}
\date{2026年10月22日初稿}
\makeatletter
\renewcommand{\maketitle}{%
  \thispagestyle{plain}%
  \begin{center}
    {\LARGE\bfseries \@title\par}
    \vspace{0.6em}
    {\small \@date\par}
  \end{center}
  \vspace{0.6em}}
\makeatother
\newcommand{\draftnotice}{\noindent{\small
本文为供逐篇审阅的学习笔记初稿，不是新增研究成果。
文中区分直接证明与引用结果；正文不依赖本文。
本轮已作公式、条件及排版初步核查，尚待作者逐句审定。}\par\medskip}

\XeTeXlinebreaklocale "zh"
\renewcommand{\proofname}{证明}
\date{2026年10月24日初稿}
\title{$1+2+3+\cdots=-\dfrac1{12}$ 的几种严格含义}
\begin{document}
\maketitle

\noindent{\small 本篇讨论发散级数 $\sum n$ 的赋值问题。凡标为命题者均给出证明；
引用的结果注明出处；数值结果单独标明，不作为证明使用。}\par\medskip

\section{问题的提法}
$\sum_{n\le N}n=N(N+1)/2\to+\infty$，故等式 $1+2+3+\cdots=-\tfrac1{12}$
不是通常意义下的级数求和。需要回答的是：在何种明确的赋值规则下该等式成立；
这些规则须放弃哪些性质；不同规则为何给出同一常数。

本文给出两种互相独立的严格表述：$\zeta$ 的解析延拓在 $s=-1$ 处的值（第~3 节），
以及光滑截断部分和在分离发散主项后的常数项（第~4 节），
并说明二者均来自 Euler--Maclaurin 公式在左端点的同一边界项；第~6 节据此辨析几种流行说法。

\section{求和法的公理性质与不相容命题}
称一个由某线性数列空间到 $\mathbb C$ 的映射 $V$ 为\emph{求和法}。考虑以下性质：
\begin{enumerate}
 \item[(R)] 正则：对收敛级数，$V$ 给出通常的和；
 \item[(M)] 单调：对实数列，若 $a_n\le b_n$ 对一切 $n$ 成立，则 $V(a)\le V(b)$；
 \item[(L)] 线性：$V(a+b)=V(a)+V(b)$，$V(ca)=cV(a)$；
 \item[(S)] 平移稳定：$V(a_1,a_2,\ldots)=a_1+V(a_2,a_3,\ldots)$。
\end{enumerate}
通常的收敛求和、Cesàro 求和与 Abel 求和同时满足以上四条。

\begin{proposition}\label{prop:monotone}
若 $V$ 满足 (R)(M)，且 $(1,2,3,\ldots)$ 在其定义域内，则 $V(1,2,3,\ldots)\ge N(N+1)/2$ 对一切 $N$ 成立。
因而不存在满足 (R)(M) 并给 $\sum n$ 赋予有限值的求和法。
\end{proposition}
\begin{proof}
数列 $(1,2,\ldots,N,0,0,\ldots)$ 对应收敛级数，由 (R) 其值为 $N(N+1)/2$；
它逐项不超过 $(1,2,3,\ldots)$，由 (M) 得结论。
\end{proof}

\begin{proposition}\label{prop:triangle}
不存在满足 (L)(S)、其定义域含 $(1,2,3,\ldots)$ 并对去掉首项封闭的复值求和法。
特别地，满足 (S) 的求和法不能给 $(1,1,1,\ldots)$ 赋值。
\end{proposition}
\begin{proof}
设 $S=V(1,2,3,\ldots)$。由 (S)，$S=1+V(2,3,4,\ldots)$。
由 (L)，$V(2,3,4,\ldots)=S+T$，其中 $T=V(1,1,1,\ldots)$，
其存在性由定义域的线性得到。于是 $T=-1$。
另一方面，对 $(1,1,1,\ldots)$ 用 (S) 得 $T=1+T$，矛盾。
\end{proof}

因此，在保留 (R) 并包含有限支集数列的前提下，有限赋值必须放弃 (M)；
在保留 (L) 且定义域对去掉首项封闭的前提下，还必须放弃一般的 (S)。下面两种规则均保留 (L)，
其值依赖于各项所处的\emph{位置} $n$，而不只依赖于各项构成的多重集。

\section{\texorpdfstring{第一种表述：$\zeta$ 的延拓在 $s=-1$ 处的值}{第一种表述：zeta 的延拓在 s=-1 处的值}}
\subsection{周期 Bernoulli 函数}
记 $t=\{x\}$，令
\[
 P_1(x)=t-\tfrac12,\qquad
 P_2(x)=\tfrac12\bigl(t^2-t+\tfrac16\bigr),\qquad
 P_3(x)=\tfrac16\bigl(t^3-\tfrac32t^2+\tfrac12t\bigr),
\]
即 $P_k(x)=B_k(\{x\})/k!$，$B_k$ 为 Bernoulli 多项式。直接验证：
$P_2,P_3$ 连续、以 $1$ 为周期，在非整数点 $P_2'=P_1$、$P_3'=P_2$，并且
\[
 P_2(0)=P_2(1)=\tfrac1{12}=\tfrac{B_2}{2},\qquad P_3(0)=P_3(1)=0,\qquad
 \max|P_3|=\tfrac{\sqrt3}{216}.
\]
最后一式在 $t=\tfrac12\pm\tfrac{\sqrt3}6$ 处取到。

\subsection{\texorpdfstring{延拓到 $\Rea s>-2$}{延拓到 Re s>-2}}
学习笔记第~4 篇证明了，对 $\Rea s>0$、$s\ne1$，
\[
 \zeta(s)=\frac{s}{s-1}-s\int_1^\infty\{x\}x^{-s-1}\dd x .
\]
将 $\{x\}=P_1(x)+\tfrac12$ 代入，并用 $s\int_1^\infty\tfrac12x^{-s-1}\dd x=\tfrac12$，得
\[
 \zeta(s)=\frac{s}{s-1}-\frac12-s\int_1^\infty P_1(x)x^{-s-1}\dd x,\qquad\Rea s>0.
\]
$\Rea s>0$ 时 $P_2(x)x^{-s-1}$ 与 $P_3(x)x^{-s-2}$ 在无穷远处趋于零，分部积分两次：
\[
 \int_1^\infty P_1x^{-s-1}\dd x
 =-P_2(1)+(s+1)\int_1^\infty P_2x^{-s-2}\dd x
 =-\frac1{12}+(s+1)(s+2)\int_1^\infty P_3x^{-s-3}\dd x .
\]
\begin{proposition}\label{prop:zeta}
对 $\Rea s>-2$、$s\ne1$，
\[
 \zeta(s)=\frac{s}{s-1}-\frac12+\frac{s}{12}
 -s(s+1)(s+2)\int_1^\infty P_3(x)\,x^{-s-3}\dd x .               \tag{3.1}
\]
特别地，$\zeta(0)=-\tfrac12$，$\zeta(-1)=-\tfrac1{12}$。
\end{proposition}
\begin{proof}
上面已在 $\Rea s>0$ 上证明 (3.1)。在 $\Rea s\ge-2+\delta$ 的紧集上，
被积函数由 $\tfrac{\sqrt3}{216}x^{-1-\delta}$ 控制，且对 $s$ 解析，
故积分在 $\Rea s>-2$ 内解析。右端在 $\Rea s>-2$、$s\ne1$ 上解析，
由解析延拓的唯一性，它就是 $\zeta$ 在该区域的延拓。
在 $s=0$ 与 $s=-1$ 处，积分项分别因因子 $s$ 与 $s+1$ 而为零，代入即得。
\end{proof}

在 (3.1) 中，$\zeta(-1)$ 的三项分别为 $\tfrac12$、$-\tfrac12$ 与 $-\tfrac1{12}$；
后者来自分部积分在左端点 $x=1$ 的边界项 $-P_2(1)=-B_2/2$。
继续分部积分可将 $\zeta$ 延拓到全平面，并得到
$\zeta(-n)=-B_{n+1}/(n+1)$（$n\ge1$）；此处引用该结果，不再展开。

\section{第二种表述：光滑截断下的常数项}
\subsection{Euler--Maclaurin 型恒等式}
\begin{lemma}\label{lem:em}
设 $f\in C^3[0,\infty)$，且在某有限区间外为零。则
\[
 \sum_{n\ge1}f(n)=\int_0^\infty f\dd x-\frac{f(0)}2-\frac{f'(0)}{12}
 +\int_0^\infty P_3(x)f'''(x)\dd x .                                \tag{4.1}
\]
\end{lemma}
\begin{proof}
在 $[n,n+1]$ 上，$\int_n^{n+1}P_1f'\dd x=\tfrac12(f(n)+f(n+1))-\int_n^{n+1}f\dd x$。
对 $n\ge0$ 求和（仅有限项非零）得
$\sum_{n\ge1}f(n)=\int_0^\infty f-\tfrac12f(0)+\int_0^\infty P_1f'$。
由 $P_2'=P_1$、$P_3'=P_2$ 及 $f$ 的紧支性，
$\int_0^\infty P_1f'=-P_2(0)f'(0)-\int_0^\infty P_2f''
=-\tfrac1{12}f'(0)+\int_0^\infty P_3f'''$，其中用了 $P_3(0)=0$。
\end{proof}

\subsection{主命题}
设 $\eta\in C^\infty[0,\infty)$，在某有限区间外为零，且 $\eta(0)=1$。令
\[
 S_\eta(N)=\sum_{n\ge1}n\,\eta(n/N),\qquad C_\eta=\int_0^\infty y\,\eta(y)\dd y .
\]
\begin{proposition}\label{prop:smooth}
当 $N\to\infty$ 时，
\[
 S_\eta(N)=C_\eta N^2-\frac1{12}+O_\eta(N^{-1}).
\]
\end{proposition}
\begin{proof}
对 $f(x)=x\,\eta(x/N)$ 用引理~\ref{lem:em}。有
$\int_0^\infty f=C_\eta N^2$，$f(0)=0$，$f'(0)=\eta(0)=1$，以及
\[
 \int_0^\infty|f'''|\dd x
 \le\frac1N\int_0^\infty\bigl(3|\eta''(y)|+y|\eta'''(y)|\bigr)\dd y .
\]
结合 $|P_3|\le\sqrt3/216$ 即得。
\end{proof}

主项 $C_\eta N^2$ 依赖截断函数，常数项 $-\tfrac1{12}$ 与 $\eta$ 无关。
由引理~\ref{lem:em} 可见，它就是 $-f'(0)/12=-B_2f'(0)/2$，
即第~3 节中同一个左端点边界项。两种表述给出同一常数，原因在此。

\subsection{与硬截断的比较}
硬截断 $\sum_{n\le N}n=\tfrac12N^2+\tfrac12N$ 的常数项为 $0$。
将引理~\ref{lem:em} 的推导限制在 $[0,N]$ 上，右端点 $x=N$ 另外贡献
$\tfrac12f(N)+\tfrac1{12}f'(N)$；对 $f(x)=x$，即 $\tfrac12N+\tfrac1{12}$。
左端点的 $-\tfrac1{12}$ 与右端点的 $+\tfrac1{12}$ 相消，同时产生 $\tfrac12N$ 项。
光滑截断函数在支集右端各阶导数为零，不产生右端点边界项，
因此常数项只保留左端点的贡献。

\subsection{数值检验与高阶修正}
取 $\eta(x)=\exp\bigl(1-\tfrac1{1-x^2}\bigr)$（$0\le x<1$），$x\ge1$ 时 $\eta(x)=0$。
它满足命题~\ref{prop:smooth} 的条件，$C_\eta=0.2018263188\ldots$。
由于 $\eta$ 在 $[1,\infty)$ 上为零，$S_\eta(N)$ 是有限和。
用 \texttt{mpmath} 在 40 位工作精度下计算，得：
\begin{center}
\begin{tabular}{rll}
\toprule
$N$ & $S_\eta(N)-C_\eta N^2$ & $\bigl(S_\eta(N)-C_\eta N^2+\tfrac1{12}\bigr)N^2$\\
\midrule
10   & $-0.0922356\ldots$ & $-0.890\ldots$\\
30   & $-0.0834221\ldots$ & $-0.0799\ldots$\\
100  & $-0.0833341\ldots$ & $-0.008303\ldots$\\
300  & $-0.0833334259\ldots$ & $-0.0083333\ldots$\\
1000 & $-0.0833333416\ldots$ & $-0.0083333313\ldots$\\
\bottomrule
\end{tabular}
\end{center}
第三列表明本例误差为 $N^{-2}$ 量级，且趋于 $-\tfrac1{120}$。这是数值观察，以下给出解释。

将 (4.1) 中的余项再分部积分两次，
利用 $P_4=B_4(\{x\})/24$、$P_5=B_5(\{x\})/120$，$P_4(0)=B_4/24=-\tfrac1{720}$，$P_5(0)=0$，得
\[
 \int_0^\infty P_3f'''\dd x=\frac{f'''(0)}{720}+\int_0^\infty P_5f^{(5)}\dd x .
\]
对 $f(x)=x\,\eta(x/N)$，$f'''(0)=3\eta''(0)/N^2$，且 $\int_0^\infty|f^{(5)}|\dd x=O_\eta(N^{-3})$。故
\[
 S_\eta(N)=C_\eta N^2-\frac1{12}+\frac{\eta''(0)}{240\,N^2}+O_\eta(N^{-3}).   \tag{4.2}
\]
本例 $\eta(x)=1-x^2+O(x^4)$，$\eta''(0)=-2$，修正项为 $-\tfrac1{120}N^{-2}$，与数值一致。
注意 $\tfrac1{240}=\tfrac12\zeta(-3)$：(4.2) 中 $N^{-2}$ 的系数是 $\zeta(-3)\eta''(0)/2!$，
常数项是 $\zeta(-1)\eta(0)/0!$。

\begin{remark}
一般地，用 Mellin 反演写
$S_\eta(N)=\frac1{2\pi\ii}\int_{(c)}\zeta(s-1)\tilde\eta(s)N^s\dd s$（$c>2$，$\tilde\eta$ 为 $\eta$ 的 Mellin 变换），
$\tilde\eta$ 在 $s=-k$ 处有留数 $\eta^{(k)}(0)/k!$，左移积分线可得
$N^{-k}$ 项的系数为 $\zeta(-k-1)\eta^{(k)}(0)/k!$。
由于 $\zeta(-2)=0$，$N^{-1}$ 项消失，与本例观察相符。
移线所需的 $\tilde\eta$ 在竖线上的快速衰减与 $\zeta$ 的多项式增长此处不再验证；
参见 Tao 的文章（文末参考）。学习笔记第~1、2 篇讨论了 Mellin 反演与光滑权重。
\end{remark}

对 $(1,1,1,\ldots)$ 作同样的计算，引理~\ref{lem:em} 给出常数项 $-\tfrac12=\zeta(0)$；
上述 $\eta$ 下 $N=100$ 时数值偏差已小于 $3\times10^{-11}$。

\section{Dirichlet 正则化下的形式运算}
\subsection{定义}
对数列 $a=(a_n)$，若 $D_a(s)=\sum a_nn^{-s}$ 在某右半平面收敛，
且可延拓为在 $s=0$ 的邻域内解析的函数，则定义 $V_D(a)=D_a(0)$。
$V_D$ 满足 (L)，不满足 (S)、(M)。它还具有\emph{伸缩不变性}：
若 $b_{mk}=a_k$、其余 $b_n=0$，则 $D_b(s)=m^{-s}D_a(s)$，故 $V_D(b)=V_D(a)$。

\subsection{常见推导的检验}
常见的形式推导是
\[
 S=1+2+3+\cdots,\quad 4S=0+4+0+8+\cdots,\quad S-4S=1-2+3-4+\cdots=\tfrac14 .
\]
在 $V_D$ 下逐步检验：
\begin{itemize}
 \item $D(s)=\zeta(s-1)$，$V_D(S)=\zeta(-1)$；
 \item $0+4+0+8+\cdots$ 由 $4S$ 伸缩得到，$D(s)=4\cdot2^{-s}\zeta(s-1)$，值为 $4\zeta(-1)$；
 \item $1-2+3-4+\cdots$ 的 $D(s)=(1-2^{2-s})\zeta(s-1)$，值为 $-3\zeta(-1)$。
 另一方面，其 Abel 和为 $\lim_{x\uparrow1}x/(1+x)^2=\tfrac14$。
\end{itemize}
由 (L)，$-3\zeta(-1)=V_D(S)-V_D(4S\text{ 的伸缩})$。上述推导的每一步在 $V_D$ 下成立，
但它预设了两件未加证明的事：$V_D(S)$ 存在，即 $\zeta$ 可延拓到 $s=-1$（命题~\ref{prop:zeta}）；
$1-2+3-4+\cdots$ 的 $V_D$ 值等于其 Abel 和（此处由两边分别计算而得）。
因此该推导只是在一个已知存在的框架内的计算，不是独立的证明。

\subsection{平移的作用}
$V_D$ 不满足 (S)。例如右移一位所得的 $0+1+2+3+\cdots$，
$D(s)=\sum(n-1)n^{-s}=\zeta(s-1)-\zeta(s)$，故
\[
 V_D(0,1,2,3,\ldots)=-\tfrac1{12}+\tfrac12=\tfrac5{12}\ne V_D(1,2,3,\ldots).
\]
这与命题~\ref{prop:triangle} 一致。导致矛盾的形式推导，通常隐含了这种平移。

\section{几种流行说法的辨析}
等式 $1+2+3+\cdots=-\tfrac1{12}$ 在数学以外流传甚广，伴随着若干不准确的解释。
以下逐条给出它们在上述框架中的确切地位。

\subsection{\texorpdfstring{“全体自然数之和等于 $-\tfrac1{12}$”}{“全体自然数之和等于 −1/12”}}
若“和”指部分和的极限，则和为 $+\infty$。
由命题~\ref{prop:monotone}，任何满足 (R)(M)、且包含所用有限支集数列的复值求和法，
都不能给它赋予有限值。
$V_D$ 不满足一般的 (M)(S)；命题~\ref{prop:smooth} 则明确提取分离主项后的常数项。
不指明规则而说“等于”，是没有确定含义的陈述。

\subsection{\texorpdfstring{“$\zeta(s)=\sum n^{-s}$，代入 $s=-1$ 即得”}{“ζ(s)=Σn⁻ˢ，代入 s=−1 即得”}}
等式 $\zeta(s)=\sum n^{-s}$ 仅在 $\Rea s>1$ 成立。解析延拓的唯一性保证延拓函数唯一，
但不保证延拓函数在收敛域外仍由原级数表示。同一结构的初等例子是
\[
 \frac1{1-z}=\sum_{n\ge0}z^n\quad(|z|<1),
\]
左端在 $z=2$ 处取值 $-1$，而右端发散；由此只能说 $-1$ 是该函数的延拓在 $z=2$ 的值，
不能说 $1+2+4+8+\cdots=-1$ 在通常意义下成立。$\zeta(-1)$ 与 $\sum n$ 的关系与此相同。

\subsection{\texorpdfstring{以 $1-1+1-\cdots=\tfrac12$ 为起点的推导}{以 1−1+1−⋯=1/2 为起点的推导}}
另一种流行的推导分三步：
\begin{enumerate}
 \item $G=1-1+1-1+\cdots=\tfrac12$（取部分和的 Cesàro 平均）；
 \item $B=1-2+3-4+\cdots$，将 $B$ 与其右移一位 $0+1-2+3-\cdots$ 逐项相加得 $G$，故 $B=\tfrac14$；
 \item $S-B=0+4+0+8+0+12+\cdots=4S$，故 $S=-\tfrac1{12}$。
\end{enumerate}
第 2 步若仅凭一般的平移稳定性处理右移，第 3 步又对 $S$ 使用线性并预设其赋值存在，
就用到了命题~\ref{prop:triangle} 已排除的全局公理组合。
不过，公理组合不相容，不等于这些\emph{具体等式}不能在某个较弱框架内同时成立。

先看普通的一阶 Cesàro 求和。$B$ 的前 $2m$ 个部分和的平均为 $0$，
前 $2m+1$ 个部分和的平均为 $(m+1)/(2m+1)$，分别趋于 $0$ 与 $1/2$。
因此 $B$ 不具有这种 Cesàro 和。Abel 求和则给出 $G=1/2$、$B=1/4$，
并允许第 2 步的右移；但 $S$ 的 Abel 值不是有限数，不能据此继续第 3 步。

在 Dirichlet 正则化 $V_D$ 中，三个值却可分别核验为
$V_D(G)=1/2$、$V_D(B)=1/4$、$V_D(S)=-1/12$。
为检查第 2 步，记 $\eta(s)=(1-2^{1-s})\zeta(s)$，
右移后的 $B$ 所对应的 Dirichlet 函数为
\[
 \sum_{n\ge2}(-1)^n(n-1)n^{-s}=\eta(s)-\eta(s-1).
\]
延拓到 $s=0$ 后，其值为 $1/2-1/4=1/4$，所以这一特定右移确实不改变赋值。
第 3 步的伸缩与线性关系已在第~5.2 节核验。
这说明三条具体等式可以在 $V_D$ 下同时成立；
但它们依赖已建立的解析延拓与各式的单独验证，不能反过来充当这些前提的独立证明。

\subsection{\texorpdfstring{“物理实验证实了 $-\tfrac1{12}$”}{“物理实验证实了 −1/12”}}
在量子场论的某些计算中（例如平行导体板间 Casimir 力的标准推导），
会出现对模式能量求和的发散级数。物理上的处理是引入依赖频率的截断或正则化，
可测量的量对应于与截断无关的有限部分；其结构与命题~\ref{prop:smooth} 相同：
发散主项依赖截断，常数项不依赖截断。
实验检验的是这一正则化程序所预言的物理量，而不是“全体自然数之和”这一数学对象；
它不构成 $\sum n$ 在任何通常意义下收敛到 $-\tfrac1{12}$ 的证据。
此处只说明逻辑地位，不展开物理推导。

\section{讨论}
\begin{enumerate}
 \item 由命题~\ref{prop:monotone}，$-\tfrac1{12}$ 不能解释为保留上述正则性与单调性的有限总和。
 命题~\ref{prop:smooth} 给出一个确切解释：它是光滑截断部分和中与截断函数无关的常数项；
 发散主项 $C_\eta N^2$ 则依赖截断方式。
 \item 本文的两种表述给出相同的值，共同来源是 Euler--Maclaurin 公式的边界项 $-B_2/2$。
 Ramanujan 求和涉及另外的定义与归一化，本文不据名称相同就断言其赋值；
 有关不同求和法可参看 Hardy, \emph{Divergent Series}。
 \item 由命题~\ref{prop:triangle} 与第~5.3 节，这些赋值依赖于项的编号。
 更准确的说法是：所赋予的值属于函数 $n\mapsto n$（连同其定义域 $\mathbb Z_{\ge1}$），
 而非属于数的多重集 $\{1,2,3,\ldots\}$。
 \item 本书的一个基本做法，是在收敛域内定义 $\zeta$ 再作解析延拓。
 $\zeta(-1)=-\tfrac1{12}$ 说明，延拓所得的函数值是确定的数学对象，
 但不再是原级数的和；临界带内的 $\zeta(s)$ 也是如此。
\end{enumerate}

\begin{samepage}
\section*{参考}
\begin{itemize}
 \item G.\ H.\ Hardy, \emph{Divergent Series}, Oxford, 1949.
 \item T.\ Tao, \emph{The Euler--Maclaurin formula, Bernoulli numbers, the zeta function, and real-variable analytic continuation}, 2010.
 \url{https://terrytao.wordpress.com/2010/04/10/the-euler-maclaurin-formula-bernoulli-numbers-the-zeta-function-and-real-variable-analytic-continuation/}
 \item 本书关于 $\zeta$ 解析延拓与函数方程的章节；学习笔记第~1、2、4 篇。
\end{itemize}
\end{samepage}
\end{document}
